\documentclass[10pt,a4paper]{article}
\usepackage[margin=0.5in]{geometry}
\usepackage{amsmath}
\usepackage{graphicx}
\usepackage{booktabs}
\usepackage{setspace}
\begin{document}

\title{\vspace{-10mm}\textbf{Quarterly Technical Report: System Scalability}}
\author{Engineering Architecture Team}
\date{October 2026}
\maketitle

\section{Executive Summary}
During the third quarter, the platform experienced a 300\% increase in traffic. This technical report details the infrastructure changes implemented to ensure zero downtime and maintain sub-100ms P99 latency. The legacy monolith was completely decoupled. All core services now operate independently, communicating via asynchronous message queues to prevent cascading failures.

\section{Architecture Diagram}
\begin{figure}[h!]
\begin{center}
\includegraphics[width=0.4\linewidth]{figure_001.png}
\end{center}
\vspace{-4mm}
\caption{System diagram showing the unidirectional flow.}
\end{figure}

\section{Performance Metrics}
\begin{center}
\begin{tabular}{llr}
\toprule
Metric & Before & After \\
\midrule
Latency & 200ms & 50ms \\
Throughput & 1k & 5k \\
\bottomrule
\end{tabular}
\end{center}

\section{Load Balancing Formula}
The asynchronous distribution of load is defined as:
\begin{equation}
i\hbar \frac{\partial \Psi}{\partial t} = \hat{H}\Psi + \alpha \nabla^2 \Psi
\end{equation}

\section{Conclusion}
The migration was completed successfully within the allocated budget. Future work will explore multi-region failover.
\end{document}
